\section{Introduction}\label{sec:introduction}

CinePal is a conversational clustering system designed for exploring movie catalogs through natural dialogue.
Rather than using static operations, the system continuously partitions the database into named, inspectable clusters that the user navigates and refines throughout the conversation.
Because there is \textbf{no intrinsic ground truth for cluster quality}, the user's satisfaction serves as the explicit objective function.

The core challenge lies in inferring user \textbf{intent}, translating it into operations, and maintaining a legible structure while recovering gracefully from errors.
To sustain this workflow, \textbf{label persistency} stabilizes cluster names across turns.
At the same time, a content-addressed snapshot tree accommodates \textbf{evolving preferences} by enabling seamless undo, reset, and branching operations.

Ultimately, CinePal is a \textbf{build-heavy} project.
Its primary contribution is the system itself, which integrates a multiagent conversational architecture, a reproducible multi-modal data pipeline, and an automated evaluation harness.
